% !PRINTED 17
\weberpage{17}
%<<CONT
isomorphic, $\mu$-$\nu$-fold isomorphic to one another. By far
the most important is the simple isomorphism, which we shall
therefore designate as isomorphism pure and simple, while a
multiple isomorphism shall always be made recognisable by an
added word $^{1)}$.

\begin{center}
\S.\,6.

\gesperrt{The composition series and the theorem of C. Jordan.}
\end{center}

A group $P$ is called \gesperrt{simple} if it has no normal
divisor besides itself and unity, \gesperrt{composite} if further
normal divisors are present.

A normal divisor which has no other normal divisor \gesperrt{above it},
which therefore is not part of another proper normal divisor
of $P$, is called \gesperrt{a greatest normal divisor}
of $P$ $^{2)}$. If $P_1$ is a greatest normal divisor of $P$, $P_2$ a
greatest normal divisor of $P_1$, $P_3$ a greatest normal divisor
of $P_2$, and so on, then the series of groups

\weq{1}{P,\, P_1,\, P_2,\, P_3 \ldots 1,}

which necessarily ends with the group 1 formed from the unit
element alone, is called a \gesperrt{composition series} of the
group $P$.

We shall denote the degrees of the groups (1) by $n, n_1, n_2, n_3 \ldots 1$,
and the quotients $n : n_1$, $n_1 : n_2$, $n_2 : n_3 \ldots$, that is,
the indices of $P_1$ with respect to $P$, of $P_2$ with respect to $P_1$, and so on,
by $\nu_1, \nu_2, \nu_3 \ldots$ The series of integers

\weq{2}{\nu_1,\, \nu_2,\, \nu_3 \ldots}

we call the \gesperrt{index series} of the group $P$.

For a given group $P$ the greatest normal divisors $P_1, P_2, P_3 \ldots$
can in general be selected in several different
ways, and accordingly the degree numbers
$n_1, n_2, n_3 \ldots$ and the indices $\nu_1, \nu_2, \nu_3 \ldots$ may also turn out
%CONT>>

$^{1)}$ Cf. C.\,\gesperrt{Jordan}, Traité des substitutions. \gesperrt{Netto}, Substitutionentheorie
. The designation ``$\mu$-fold'' is due to \gesperrt{Netto}. According to \gesperrt{Jordan}
the simple isomorphism is called ``holohedral'', the multiply
``merohedral'' isomorphism. The simple isomorphism is sometimes
 also designated as equivalence.

$^{2)}$ Distinguished maximal subgroup, according to another designation.

%% [page furniture, not body text -- commented out 2026-08-10] Page furniture (running
%% head, folio, signature) is excluded from this edition by standing policy.
%% Weber, Algebra. II. \hfill 2
